\documentclass[10pt]{article}

\usepackage[margin=0.75in]{geometry}
\usepackage{amsmath,amsthm,amssymb}
\usepackage{xcolor}
\usepackage{cancel}
\usepackage{graphicx}
\usepackage{changepage}
\usepackage{circuitikz}
\usepackage{pgfplots}
\usepackage{physics}
\usepackage{hyperref}
\usepackage{siunitx}
\usepackage{fontspec}
\usepackage{relsize}
\usepackage{subfig}
\usepackage{polynom}
% \usepackage{polydiv}
\usepackage{todonotes}
\usepackage{multicol, multirow, booktabs}
\usepackage[breakable]{tcolorbox}
\usepackage[inline]{enumitem}

\theoremstyle{definition}
\newtheorem{problem}{Problem}
\newtheorem{soln}{Solution}
\newtheorem{theorem}{Theorem}

\pgfplotsset{compat=newest}
\usetikzlibrary{lindenmayersystems}
\usetikzlibrary{arrows}
\usetikzlibrary{calc}
\usetikzlibrary{positioning, fit}
\usetikzlibrary{3d, perspective}

\definecolor{incolor}{HTML}{303F9F}
\definecolor{outcolor}{HTML}{D84315}
\definecolor{cellborder}{HTML}{CFCFCF}
\definecolor{cellbackground}{HTML}{F7F7F7}
\newcommand{\ui}{\hat{i}}
\newcommand{\uj}{\hat{j}}
\newcommand{\uk}{\hat{k}}
\newcommand{\ux}{\hat{x}}
\newcommand{\uy}{\hat{y}}
\newcommand{\uz}{\hat{z}}
\newcommand{\pr}[1]{{#1}^\prime}
\newcommand{\hull}[1]{\langle {#1}\rangle}
\pgfdeclarelayer{background}  
\pgfsetlayers{background,main}
\AtBeginDocument{\RenewCommandCopy\qty\SI}
\newcommand{\justif}[2]{&{#1}&\text{#2}}

\makeatletter
\newcommand{\boxspacing}{\kern\kvtcb@left@rule\kern\kvtcb@boxsep}
\makeatother
\newcommand{\prompt}[4]{
    \ttfamily\llap{{\color{#2}[#3]:\hspace{3pt}#4}}\vspace{-\baselineskip}
}


\newcommand{\thevenin}[2]{
  \begin{center}
    \begin{circuitikz} \draw
      (0,0) -- (2,0) to[battery1, l_=$V_{Th}\eq#1$] (2,2) 
      to[resistor, l_=$R_{Th}\eq#2$] (0,2)
      ;
      \draw [o-] (-.07,2.079);
      \draw [o-] (-.07,0.079);
    \end{circuitikz}
  \end{center}
}

\newcommand{\norton}[2]{
  \begin{center}
    \begin{circuitikz} \draw
      (0,0) -- (3,0) to[american current source, l_=$I_{N}\eq#1$] (3,2) -- (0,2) (2,0)
      to[resistor, l=$R_{N}\eq#2$] (2,2)
      ;
      \draw [o-] (-.07,2.079);
      \draw [o-] (-.07,0.079);
    \end{circuitikz}
  \end{center}
}

\newcommand{\highlight}[1]{\colorbox{yellow}{$\displaystyle #1$}}

\newcommand{\ti}[1]{\widetilde{#1}}

\newfontface{\Kaufmann}{Kaufmann}
\DeclareTextFontCommand{\kf}{\Kaufmann}
\newcommand{\scriptr}{\fontsize{12pt}{12pt}\kf{r}}

\newfontface{\KaufmannB}{Kaufmann Bd BT}
\DeclareTextFontCommand{\kfb}{\KaufmannB}
\newcommand{\bscriptr}{\fontsize{12pt}{12pt}\kfb{r}}

\newcommand{\bv}[1]{\mathbf{#1}}

\title{Math 4360H: Assignment IV}
\author{Jeremy Favro (0805980) \\ Trent University, Peterborough, ON, Canada}
\date{\today}

\begin{document}
\maketitle

% PROBLEM 1
\begin{problem}
Find the degree of the extension $\mathbb{Q}(1+\sqrt[3]{2}+\sqrt[3]{4})$ over $\mathbb{Q}$.
\end{problem}
\begin{soln}
  The degree of such an extension is given by the degree of its minimal polynomial. We can find the minimal polynomial by taking
  $1+\sqrt[3]{2}+\sqrt[3]{4}=x$ and trying to solve for rational coefficients such that we obtain a polynomial in $\mathbb{Q}[x]$ which equals zero:
  \begin{align*}
             & x=1+\sqrt[3]{2}+\sqrt[3]{4}         \\
    \implies & (x-1)^3=2(1+\sqrt[3]{2})^3          \\
    \implies & (x-1)^3=6\sqrt[3]{2}+6\sqrt[3]{4}+6 \\
    \implies & (x-1)^3=6x                          \\
    \implies & x^3-3x^2+3x-1-6x=0                  \\
    \implies & x^3-3x^2-3x-1=0.                     
  \end{align*}
  Since by RZT the only possible rational zeroes of this are $\pm 1$ and neither is a zero this is irreducible and hence
  the minimal polynomial. This means $[\mathbb{Q}(1+\sqrt[3]{2}+\sqrt[3]{4}):\mathbb{Q}]=3$.
\end{soln}

% PROBLEM 2
\begin{problem}
Find the degree of the extension $\mathbb{Q}(\sqrt[3]{2},\sqrt[3]{6},\sqrt[3]{24})$ over $\mathbb{Q}$.
\end{problem}
\begin{soln}
To do this we break down the extension by tower theorem into 
$$[\mathbb{Q}(\sqrt[3]{2},\sqrt[3]{6},\sqrt[3]{24}):\mathbb{Q}]
=[\mathbb{Q}(\sqrt[3]{2})(\sqrt[3]{6})(\sqrt[3]{24}):\mathbb{Q}]
=[\mathbb{Q}(\sqrt[3]{2})(\sqrt[3]{24})(\sqrt[3]{6}):\mathbb{Q}(\sqrt[3]{24})(\sqrt[3]{6})][\mathbb{Q}(\sqrt[3]{24})(\sqrt[3]{6}):\mathbb{Q}(\sqrt[3]{24})][\mathbb{Q}(\sqrt[3]{24}):\mathbb{Q}].$$
Immediately we can see that 
$[\mathbb{Q}(\sqrt[3]{24}):\mathbb{Q}]=3$ as 24 is a root of $x^3-24$. Now we seek the minimal polynomial of $\mathbb{Q}(\sqrt[3]{24})(\sqrt[3]{6})$ over $\mathbb{Q}(\sqrt[3]{24})$.
\end{soln}

% PROBLEM 3
\begin{problem}
Let $E$ be an extension field of a field $F$ and $[E : F ]$ be a prime
number. For any $\alpha \in E\setminus F $, show that $E = F (\alpha)$.
\end{problem}
\begin{soln}
  Since $[E : F ]=p$ is prime and $\alpha \in E\setminus F\implies F\subseteq F(\alpha) \subseteq E$ we can form a tower
  which gives us that
  $[E : F ]=[E : F(\alpha) ][F(\alpha): F ]=p$.
  The only way this can happen is if $[E : F(\alpha) ]=1$ and $[F(\alpha): F ]=p$ or vice-versa. Since we are given that $\alpha \notin F$, clearly
  $[F(\alpha): F ]\neq 1 p$ so $[E : F(\alpha) ]=1\implies E=F(\alpha)$.
\end{soln}

% PROBLEM 4
\begin{problem}
Let $E$ be an extension field of a field $F$ and let $\alpha\in E$ be algebraic
of odd degree over $F$. Show that $\alpha^2$ is algebraic of odd degree over
$F$ and $F(\alpha)=F(\alpha^2)$.
\end{problem}
\begin{soln}
  $\alpha$ odd over $F$ means that $[F(\alpha):F]$ is odd. Now if $F(\alpha)\neq F(\alpha^2)$ then $\alpha \notin F(\alpha^2)$.
  We can hence construct the tower $[F(\alpha):F]=[F(\alpha):F(\alpha^2)][F(\alpha^2):F]$ which we know to be odd. Since
  $[F(\alpha):F(\alpha^2)]=2$ as $\alpha$ is a root of $x^2-\alpha^2$ the minimal polynomial for $\alpha$ over $F(\alpha^2)[x]$. Since the
  entire tower must be odd this cannot be the minimal polynomial and hence $[F(\alpha):F(\alpha^2)]=1$ so $F(\alpha)=F(\alpha^2)$. Then we can also
  see that $\alpha^2$ has odd degree as $[F(\alpha):F]=[F(\alpha):F(\alpha^2)][F(\alpha^2):F]=[F(\alpha^2):F]$ and $[F(\alpha):F]$ is odd.
\end{soln}

% PROBLEM 4
\begin{problem}
Let $E$ be a finite extension of a field $F$, and let $p(x) \in F [x]$ be an
irreducible over $F$. Show that if $\deg(p(x))$ does not divide $[E : F ]$,
then $p(x)$ has no zeros in $E$.
\end{problem}
\begin{soln}
  Suppose that $\deg(p(x))$ does not divide $[E:F]$ and that there is an $\alpha \in E$ such that $p(\alpha)=0$. This means that $[F(\alpha):F]=\deg(p(x))$
  which we can deconstruct into a tower for
  $[E:F]=[E:F(\alpha)][F(\alpha):F]=[E:F(\alpha)]\deg(p(x))$ so $\deg(p(x))$ divides $[E:F]$, a contradiction. This means our initial supposition cannot hold
  and so there cannot exist and $\alpha\in E$ for which $p(\alpha)=0$ when $\deg(p(x))$ does not divide $[E:F]$.
\end{soln}
\end{document}